% Chương 1: Giới thiệu (bản rút gọn)
\chapter{GIỚI THIỆU}
\label{chap:gioi-thieu}

\section{Bối cảnh và động lực}
\label{sec:boi-canh}

\subsection{Chi phí của một lần kiểm chứng văn bằng}

Văn bằng đại học là lời xác nhận của cơ sở đào tạo về năng lực của người học. Giá trị của lời xác nhận ấy chỉ tồn tại khi người khác kiểm chứng được nó. Trong thời kỳ thị trường lao động còn bó hẹp trong phạm vi địa phương, việc kiểm chứng diễn ra thưa thớt và được giải quyết bằng thủ tục hành chính: gửi công văn hỏi, nộp bản sao công chứng, đối chiếu với sổ gốc. Quy mô nhỏ cho phép cách làm thủ công đó vận hành ổn định.

Điều kiện này đã thay đổi. Tuyển dụng chuyển dần lên các nền tảng trực tuyến và hình thức làm việc từ xa trở nên phổ biến, nên hồ sơ ứng viên ngày càng đi xa khỏi nơi đã cấp văn bằng. Cùng lúc, số người học và số văn bằng tăng nhanh. Tần suất kiểm chứng vì thế tăng mạnh, trong khi năng lực xử lý của bộ phận lưu trữ ở từng trường hầu như giữ nguyên.

Số liệu của cơ quan quản lý cho thấy quy mô đó không còn nhỏ. Từ tháng 01/2017 đến hết tháng 11/2023, Cục Quản lý chất lượng thuộc Bộ Giáo dục và Đào tạo đã tiếp nhận 37.436 hồ sơ đề nghị công nhận văn bằng do cơ sở nước ngoài cấp và công nhận 35.662 hồ sơ; khoảng 4,7\% số hồ sơ không được công nhận. Đây mới chỉ là một luồng xác minh hẹp, vậy mà mỗi năm đã phát sinh trung bình hơn 5.000 hồ sơ. Theo Thông tư 13/2021/TT-BGDĐT \cite{TT13_2021}, thời hạn giải quyết là 20 ngày làm việc và có thể kéo dài tới 45 ngày nếu phải xác minh với cơ sở nước ngoài. Phí xác minh theo Thông tư 164/2016/TT-BTC \cite{TT164_2016} là 250.000 đồng hoặc 500.000 đồng cho mỗi văn bằng tuỳ trường hợp, chưa kể chi phí dịch thuật và công chứng.

Thời gian tính bằng tuần và chi phí tính theo từng lượt không đơn thuần là hệ quả của sự chậm trễ hành chính. Chúng xuất phát từ một đặc điểm kiến trúc: thông tin xác thực không nằm trong văn bằng mà nằm ở cơ sở đã cấp nó. Bên xác minh buộc phải hỏi lại nơi cấp, nên tổng chi phí tăng tỉ lệ thuận với số lần hỏi. Rút ngắn thời gian trả lời chỉ làm giảm hệ số của quan hệ tỉ lệ ấy chứ không xoá được nó.

\subsection{Hành lang pháp lý đã sẵn sàng}

Về chính sách, Việt Nam đã có tương đối đủ cơ sở pháp lý cho việc số hoá văn bằng. Quyết định 749/QĐ-TTg năm 2020 \cite{QD749_2020} đưa giáo dục vào nhóm lĩnh vực ưu tiên của Chương trình Chuyển đổi số quốc gia. Quyết định 131/QĐ-TTg năm 2022 \cite{QD131_2022} cụ thể hoá định hướng này cho giai đoạn 2022--2025. Luật Giao dịch điện tử số 20/2023/QH15 \cite{LuatGDDT2023} công nhận giá trị pháp lý của thông điệp dữ liệu và chữ ký điện tử.

Mốc quan trọng nhất với đề tài là Thông tư 10/2026/TT-BGDĐT, có hiệu lực từ 15/04/2026 \cite{TT10_2026}. Lần đầu tiên văn bằng số và chứng chỉ số được định danh chính thức và có giá trị pháp lý như văn bằng giấy. Thông tin trích xuất từ cơ sở dữ liệu văn bằng có giá trị như thông tin trong sổ gốc, và việc tra cứu trên nền tảng định danh điện tử có thể thay cho bản sao từ sổ gốc. Rào cản pháp lý, vốn hay được viện dẫn để trì hoãn văn bằng số, về cơ bản đã được gỡ bỏ.

Song song với đó là một ràng buộc mới. Luật Bảo vệ dữ liệu cá nhân số 91/2025/QH15, có hiệu lực từ 01/01/2026 \cite{LuatBVDLCN2025}, quy định trách nhiệm cụ thể khi thu thập, lưu trữ và chia sẻ dữ liệu cá nhân. Với hệ thống văn bằng số, quy định này biến một khuyến nghị thành ràng buộc kiến trúc: mọi thiết kế đặt dữ liệu học tập lên một hạ tầng công khai và bất biến đều phải được xem xét lại.

Về quy mô, năm 2024 cả nước có 243 cơ sở giáo dục đại học với 2.355.711 sinh viên \cite{MOET_Stats2025}, và Quyết định 452/QĐ-TTg \cite{QD452_2025} đặt mục tiêu trên 3 triệu người học vào năm 2030. Như vậy chính sách đã cho phép văn bằng số, dữ liệu đã đủ lớn để đòi hỏi tự động hoá, nhưng cơ chế xác minh ở phần lớn cơ sở đào tạo vẫn dựa vào công văn và tra cứu nội bộ. Khoảng cách giữa khung pháp lý và năng lực kỹ thuật là không gian mà khoá luận hướng tới.

\section{Giới hạn kiến trúc của mô hình xác minh hiện hành}
\label{sec:gioi-han-mo-hinh}

Có hai loại hạn chế cần tách bạch. Loại thứ nhất do triển khai chưa đầy đủ và có thể khắc phục bằng đầu tư thêm. Loại thứ hai nằm trong cấu trúc của chính mô hình và không biến mất dù đầu tư bao nhiêu.

Việt Nam đã có các nền tảng dữ liệu tập trung như hệ thống cơ sở dữ liệu giáo dục đại học HEMIS và hệ thống tra cứu văn bằng quốc gia. Mức độ triển khai giữa các đơn vị chưa đồng đều, nhưng đó là hạn chế loại thứ nhất. Hạn chế loại thứ hai nằm ở chỗ khác: trong mô hình tập trung, bên xác minh thực ra không tin vào dữ liệu mà tin vào tổ chức vận hành cơ sở dữ liệu. Một bản ghi trả về từ cổng tra cứu chỉ đáng tin trong chừng mực người ta tin rằng cổng không bị can thiệp và tổ chức vận hành vẫn còn tồn tại. Từ đặc điểm này kéo theo ba hệ quả mà số hoá thêm không giải quyết được:

\begin{itemize}
    \item \textbf{Không kiểm chứng được một cách độc lập.} Khi bên vận hành không phản hồi, chẳng hạn vì ngoài giờ làm việc, hồ sơ thất lạc hay đơn vị cấp đã sáp nhập, một văn bằng hoàn toàn hợp lệ cũng trở nên không thể kiểm chứng.
    \item \textbf{Không liên thông ở mức định danh.} Mỗi cơ sở cấp cho người học một mã riêng; người học qua nhiều chương trình ở nhiều đơn vị sẽ có nhiều hồ sơ rời rạc, không gộp được thành một bằng chứng thống nhất.
    \item \textbf{Không biểu diễn được thẩm quyền.} Cơ sở dữ liệu trả lời được câu hỏi ``bản ghi này có tồn tại không'', nhưng không trả lời được câu hỏi ``đơn vị tạo ra bản ghi có quyền cấp loại văn bằng này không''. Trong một hệ thống nhiều cơ sở, câu hỏi thứ hai mới là câu hỏi quyết định.
\end{itemize}

Vụ việc ở Trường Đại học Đông Đô minh hoạ rõ hệ quả thứ ba. Hàng trăm văn bằng được cấp trái phép, một số đã nằm trong hồ sơ cán bộ, công chức. Những văn bằng này tồn tại thật trong hệ thống của nơi cấp và qua được mọi phép kiểm tra kiểu ``bản ghi có tồn tại không''. Lỗ hổng không nằm ở tính toàn vẹn dữ liệu mà ở việc không có cơ chế nào ràng buộc hành vi cấp phát với phạm vi thẩm quyền đã được công nhận.

Ở quy mô quốc tế, Ezell và Bear mô tả hoạt động bán văn bằng không qua đào tạo như một ngành kinh doanh trị giá hàng tỷ đô la với hơn một triệu văn bằng đã được bán \cite{EzellBear2012}. Dịch vụ xác minh văn bằng của Vương quốc Anh đã xử lý hơn hai triệu yêu cầu từ năm 2011 và phát hiện hơn 300 cơ sở mạo danh đại học, nhưng chỉ khoảng 20\% nhà tuyển dụng thực sự đối chiếu bằng cấp với nơi cấp \cite{Prospects_Hedd2026}. Con số sau cùng cho thấy bản chất của vấn đề. Gian lận tồn tại được không phải vì không phát hiện nổi mà vì chi phí phát hiện cao hơn lợi ích của từng lần kiểm tra riêng lẻ. Bài toán vì thế nên được phát biểu lại: cần hạ chi phí của một lần kiểm chứng xuống đủ thấp để xác minh trở thành thói quen mặc định. Muốn vậy thì phải kiểm chứng được mà không cần hỏi ai, và đó là yêu cầu về kiến trúc chứ không phải về quy trình.

\subsection{Hướng tiếp cận dựa trên bằng chứng mật mã}

Trong khoảng mười năm trở lại đây, các nghiên cứu và triển khai cho thấy một sự chuyển dịch khá nhất quán: thay vì số hoá tài liệu rồi tìm cách bảo vệ nó, người ta số hoá chính bằng chứng về tính xác thực. Chữ ký số gắn với nội dung cho phép phát hiện mọi thay đổi; hàm băm mật mã cho phép đối chiếu nội dung với một cam kết đã công bố; sổ cái phân tán cho phép công bố cam kết đó theo cách không sửa lại được. Kết hợp ba thành phần này, bên xác minh có thể tự kiểm tra bằng toán học thay vì phải tin vào một tổ chức.

W3C đã chuẩn hoá hai khối nền cho hướng tiếp cận này: định danh phi tập trung (DID) \cite{W3C_DIDCore2022} và mô hình dữ liệu chứng thực có thể xác minh (VC), đạt trạng thái Recommendation ở phiên bản 2.0 vào tháng 5/2025 \cite{W3C_VCDM2025}. Hai đặc tả này là nền chung của nhiều hệ thống, trong đó có hạ tầng văn bằng số của Liên minh Châu Âu \cite{EDC_ELM} và dự án Blockcerts \cite{Blockcerts2016}.

Tuy vậy, chuẩn hoá ở tầng dữ liệu không giải quyết trọn bài toán. Chữ ký hợp lệ chứng minh dữ liệu chưa bị sửa và được tạo bởi người giữ một khoá cụ thể. Nó không chứng minh được người giữ khoá có thẩm quyền cấp loại văn bằng đó. Khi ai cũng có thể sinh khoá, tạo DID và phát hành chứng thực hợp lệ về mặt kỹ thuật, câu hỏi ``nên tin bên phát hành nào'' vẫn còn nguyên. Rà soát 63 công trình giai đoạn 2018--2023 về quản trị trong kiểm định giáo dục trên nền blockchain, De Alwis và cộng sự nhận thấy hầu như không công trình nào đề xuất một khung quản trị \cite{DeAlwis2025}; các khảo sát về định danh tự chủ cũng xếp việc thiết lập niềm tin ở mức hệ sinh thái vào nhóm vấn đề ít được giải quyết nhất \cite{Schardong2022, Sedlmeir2021}. Khoảng trống giữa niềm tin mật mã và niềm tin thể chế là điểm xuất phát của khoá luận.

\section{Bài toán nghiên cứu}
\label{sec:bai-toan}

\subsection{Trường hợp nghiên cứu: Đại học Quốc gia Hà Nội}

Khoá luận chọn Đại học Quốc gia Hà Nội (ĐHQGHN) làm trường hợp nghiên cứu, dựa trên cấu trúc hơn là quy mô. ĐHQGHN là một pháp nhân duy nhất gồm 9 trường đại học thành viên và 3 trường trực thuộc, với 483 chương trình đào tạo \cite{VNU_Report2026}. Mười hai đơn vị này vận hành tương đối độc lập về quản lý nhưng cùng cấp văn bằng dưới thẩm quyền dẫn xuất từ một thực thể cấp trên. ĐHQGHN còn khuyến khích học song ngành, chuyển đổi tín chỉ và liên kết giữa các trường, nên một người học thường có quan hệ học thuật với nhiều đơn vị. Nói cách khác, cả vấn đề chuỗi thẩm quyền lẫn vấn đề liên thông định danh đều hiện diện trong một phạm vi quan sát được.

ĐHQGHN được dùng làm bối cảnh thiết kế và kiểm chứng, không phải giới hạn của mô hình. Cấu trúc ``một thực thể gốc uỷ quyền cho nhiều đơn vị cấp'' cũng là cấu trúc của hệ thống giáo dục ở cấp quốc gia, nơi Bộ Giáo dục và Đào tạo giữ vai trò tương tự ĐHQGHN trên phạm vi rộng hơn.

\subsection{Ba vấn đề cốt lõi}

Bài toán được tách thành ba vấn đề ứng với ba tầng của hệ thống. Mỗi vấn đề cần một loại giải pháp khác nhau, và giải quyết vấn đề này không tự động giải quyết vấn đề kia.

\textit{Vấn đề thứ nhất: xác minh phụ thuộc vào sự hợp tác của nơi cấp.} Muốn xác minh một văn bằng của trường thành viên, nhà tuyển dụng phải gửi yêu cầu tới phòng đào tạo và chờ trả lời. Văn bằng có giá trị suốt đời người học, còn một đơn vị hành chính thì không tồn tại mãi. Cần chuyển thông tin xác thực từ nơi cấp sang chính văn bằng và một hạ tầng công bố mà ai cũng truy cập được.

\textit{Vấn đề thứ hai: thiếu định danh liên thông cho người học.} Sinh viên học song song hai trường có hai mã sinh viên và hai hồ sơ độc lập. Mã sinh viên chỉ có nghĩa trong hệ thống đã cấp ra nó. Nghiên cứu về nền tảng tín chỉ liên trường EduCTX cho thấy đây là vấn đề chung của môi trường giáo dục đa cơ sở \cite{Turkanovic2018}.

\textit{Vấn đề thứ ba: không có cơ chế biểu diễn và kiểm chứng thẩm quyền cấp phát.} Giả sử ĐHQGHN đã có hệ thống ký số văn bằng nội bộ. Bên xác minh kiểm được chữ ký, nhưng vẫn chưa biết tổ chức đã ký có được phép cấp loại văn bằng đó hay không, và sự cho phép ấy còn hiệu lực vào thời điểm cấp hay không. Trả lời câu hỏi này đòi hỏi một lớp quản trị nằm trên lớp kỹ thuật, cho phép thực thể gốc uỷ quyền, giới hạn phạm vi, thu hồi, và quan trọng nhất là cho phép bên thứ ba tự kiểm chứng trạng thái uỷ quyền.

Ba vấn đề phụ thuộc lẫn nhau: không có định danh bền vững thì không gắn được thẩm quyền vào một thực thể cụ thể; không kiểm chứng được thẩm quyền thì xác minh độc lập chỉ xác nhận được tính toàn vẹn chứ chưa xác nhận được tính hợp lệ. Bài toán vì vậy được phát biểu như sau:

\begin{quote}
\textit{Xây dựng một kiến trúc văn bằng số cho môi trường giáo dục đa tổ chức, trong đó thông tin học thuật được xác minh độc lập bằng bằng chứng mật mã mà không cần liên hệ cơ sở cấp; người học có một định danh bền vững dùng chung giữa các cơ sở; và thẩm quyền cấp phát của mỗi tổ chức được thiết lập, giới hạn, thu hồi và kiểm chứng qua một cơ chế quản trị tra cứu công khai được.}
\end{quote}

\section{Câu hỏi, mục tiêu và phạm vi nghiên cứu}
\label{sec:cau-hoi}

Bài toán được triển khai thành ba câu hỏi nghiên cứu:

\begin{itemize}
    \item \textbf{CH1.} Kiến trúc nào cho phép cấp, lưu trữ, trình bày và xác minh văn bằng số theo chuẩn mở, sao cho bên xác minh kiểm chứng được nguồn gốc, tính toàn vẹn và trạng thái hiệu lực mà không cần tương tác với cơ sở cấp?
    \item \textbf{CH2.} Khung quản trị nào biểu diễn được chuỗi thẩm quyền trong hệ sinh thái nhiều cơ sở, gồm uỷ quyền có phạm vi, uỷ quyền nhiều cấp, giới hạn thời hạn và thu hồi, sao cho mọi thay đổi đều để lại dấu vết và tra cứu độc lập được?
    \item \textbf{CH3.} Mô hình đề xuất có khả thi về kỹ thuật không, xét trên tính đúng đắn của logic phân quyền, khả năng phát hiện gian lận, chi phí của các phép mật mã và hiệu năng của đường xác minh, kể cả khi dữ liệu tăng lên?
\end{itemize}

Ba mục tiêu tương ứng được xếp theo quan hệ phụ thuộc, mỗi mục tiêu kèm một tiêu chí hoàn thành quan sát được. \textit{Mục tiêu 1} là thiết kế kiến trúc hỗ trợ trọn vòng đời văn bằng theo chuẩn W3C trên nền một hạ tầng neo trạng thái bất biến; tiêu chí hoàn thành là hệ thống chạy được từ đăng ký định danh tới xác minh và thu hồi, và quy trình xác minh báo rõ bước nào không đạt. \textit{Mục tiêu 2} là đề xuất khung quản trị niềm tin; tiêu chí là logic phân quyền được viết thành mã thi hành và mỗi quy tắc có ít nhất một phép kiểm chứng minh nó thực sự chặn được hành vi vi phạm. \textit{Mục tiêu 3} là hiện thực hoá và đánh giá; tiêu chí là các kịch bản chạy thành công đầu cuối, các kịch bản gian lận định nghĩa trước đều bị phát hiện đúng nguyên nhân, và kết quả đo được thu bằng phương pháp mô tả rõ, lặp lại được.

Phạm vi được giới hạn theo ba chiều. Về lĩnh vực, đối tượng là văn bằng và chứng chỉ trong giáo dục đại học. Về công nghệ, khoá luận dùng các chuẩn mở về định danh phi tập trung và chứng thực có thể xác minh cùng các nguyên thuỷ mật mã đã chuẩn hoá; tiết lộ chọn lọc ở mức thuộc tính và chứng minh không tiết lộ thông tin không được triển khai. Về mức độ, hệ thống dừng ở mức kiểm chứng khái niệm: không tích hợp hệ thống quản lý đào tạo thật, không triển khai trên mạng chính và không có nghiên cứu với người dùng thật.

\section{Phương pháp nghiên cứu}
\label{sec:phuong-phap}

Khoá luận tạo ra một hiện vật, gồm một kiến trúc và một hệ thống hiện thực hoá nó, nhằm giải quyết một vấn đề thực tiễn. Dạng nghiên cứu này thuộc khoa học thiết kế, nơi tri thức hình thành qua việc thiết kế và đánh giá hiện vật \cite{Hevner2004}. Quy trình đi theo sáu bước của Peffers và cộng sự \cite{Peffers2007}: xác định vấn đề, xác lập mục tiêu của giải pháp, thiết kế và phát triển, trình diễn, đánh giá, và truyền đạt. Bước một và hai tương ứng với Chương \ref{chap:gioi-thieu} và Chương \ref{chap:co-so}; bước ba là Chương \ref{chap:kien-truc}; bước bốn là Chương \ref{chap:hien-thuc}; bước năm là Chương \ref{chap:danh-gia}; bước sáu là Chương \ref{chap:ket-luan}. Khung này đòi hỏi tiêu chí đánh giá phải được đặt \emph{trước} khi thiết kế, và các tiêu chí hoàn thành ở Mục \ref{sec:cau-hoi} cùng bộ yêu cầu ở Mục \ref{sec:yeu-cau} giữ đúng vai trò đó.

Phần đánh giá dựa trên một nguyên tắc được áp dụng nhất quán: phân biệt phép kiểm chứng minh được một tính chất với phép kiểm chỉ cho thấy hệ thống không báo lỗi. Một bước kiểm tra luôn cho kết quả đạt còn nguy hiểm hơn việc không có bước đó, vì nó tạo cảm giác an toàn không có cơ sở. Do đó mỗi rào chắn được thêm vào đều đi kèm một phép kiểm được viết để làm nó thất bại. Nguyên tắc này cũng được áp dụng cho chính các phép đo hiệu năng: một con số chỉ được giữ lại khi đã kiểm tra rằng chạy lại cho kết quả tương tự.

So với bản báo cáo trước, phần đánh giá được mở rộng thành sáu thí nghiệm có mục đích riêng. Ba thí nghiệm đo các phép thuần tính toán (chi phí mật mã phía máy khách, khả năng khôi phục nội dung từ giá trị băm trên chuỗi, kích thước danh sách trạng thái), một thí nghiệm kiểm tra giả thuyết về nguồn nhiễu của phép đo cũ, một thí nghiệm đo hành vi khi dữ liệu tăng hai bậc độ lớn, và một thí nghiệm đo chi phí thật phía chuỗi. Mỗi thí nghiệm được thiết kế để trả lời một câu hỏi mà bản trước để ngỏ.

\section{Đóng góp và cấu trúc báo cáo}
\label{sec:dong-gop}

Các đóng góp được trình bày kèm mức độ, vì một đóng góp mang tính tích hợp có giá trị khác với một cơ chế mới.

\textit{Đóng góp mới.} Thứ nhất là mô hình uỷ quyền nhiều cấp với ba ràng buộc được hợp đồng thi hành: phạm vi của cấp dưới là tập con của cấp trên, thời hạn không vượt thời hạn của cấp trên, và quyền uỷ quyền không truyền tiếp. Đặc tả giao thức truy vấn sổ đăng ký tin cậy hiện còn hoãn loại truy vấn này sang phiên bản sau \cite{ToIP_TRQP2026}. Thứ hai là hai nguyên tắc thiết kế rút ra từ việc dựng kịch bản thất bại: thao tác thu hồi phải có điều kiện lỏng hơn thao tác cấp, và một thao tác không thể hoàn tác phải được bên bị ảnh hưởng chứng minh điều kiện bằng hành động của chính họ. ``Mới'' ở đây nghĩa là không tìm thấy trong phạm vi tài liệu đã khảo sát \cite{Krul2024, Turco2026, Ambati2026}, không phải khẳng định tuyệt đối.

\textit{Đóng góp tích hợp.} Một kiến trúc neo thẩm quyền và trạng thái chứng thực trên cùng một hạ tầng, kèm quy trình xác minh bảy bước báo cáo theo từng bước, trong đó có bước đối chiếu nội dung với giá trị băm đã neo.

\textit{Đóng góp kiểm chứng.} Một bộ chứng cứ phủ định cho logic phân quyền, trong đó phần lớn phép kiểm dựng hành vi vi phạm và khẳng định đúng mã lỗi; cùng bộ sáu thí nghiệm định lượng ở Chương \ref{chap:danh-gia}. Hai kết quả của bộ thí nghiệm này có giá trị vượt ra ngoài hệ thống cụ thể: giá trị băm không muối trên chuỗi khôi phục được thuộc tính văn bằng trong vài chục mili giây khi kẻ tấn công biết các trường còn lại, và một truy vấn thiếu chỉ mục làm thời gian xác minh tăng tuyến tính theo kích thước dữ liệu.

\textit{Đóng góp bối cảnh.} Ánh xạ mô hình vào giáo dục đại học Việt Nam đúng lúc ba điều kiện vừa hội đủ: văn bằng số có giá trị pháp lý \cite{TT10_2026}, quy định bảo vệ dữ liệu cá nhân có hiệu lực \cite{LuatBVDLCN2025}, và hạ tầng số quốc gia đang hình thành \cite{NDAChain2025}.

Khoá luận \emph{không} đề xuất cơ chế mật mã mới, \emph{không} chứng minh hình thức tính đúng đắn và \emph{không} đánh giá mức độ chấp nhận của người dùng; hai câu hỏi sau cùng không suy ra được từ tính khả thi kỹ thuật \cite{SSI_Adoption2025}.

Phần còn lại của báo cáo gồm năm chương. Chương \ref{chap:co-so} trình bày cơ sở lý thuyết, khảo sát các hệ thống hiện có dưới lăng kính ba tầng niềm tin và phân loại khoảng trống nghiên cứu. Chương \ref{chap:kien-truc} đề xuất kiến trúc và khung quản trị niềm tin. Chương \ref{chap:hien-thuc} mô tả việc hiện thực hoá. Chương \ref{chap:danh-gia} trình bày sáu thí nghiệm đánh giá cùng các mối đe doạ đối với tính hợp lệ của kết quả. Chương \ref{chap:ket-luan} trả lời ba câu hỏi nghiên cứu, nêu hạn chế và hướng phát triển.
